\begin{frame}{과대적합(overfitting)의 교훈}
  \small
  \begin{center}
  \includegraphics[width=0.68\linewidth]{ch01_overfit_schematic.png}
  \end{center}
  \vskip0.05em
  \footnotesize
  \begin{itemize}
    \item 직선(과소적합)은 데이터의 굴곡을 놓치고, 16차 다항식(과대적합)은
      \textbf{점과 점 사이의 잡음까지} 외워 양 끝에서 요동친다 ---
      실데이터로는 degree 1 vs 24 다항식 회귀에서 훈련
      $R^2=1.0000$(완벽해 보임)인데 테스트 $R^2$는 \textbf{마이너스
      6만}(``평균만 예측''보다 몇만 배 나쁨).
    \item 25개 파라미터가 25개 점에 정확히 맞춰지는 순간, 다항식은
      점 사이에 있던 \textbf{잡음 자체}를 외워버림.
  \end{itemize}
  \vskip0.2em
  \begin{block}{핵심 교훈}
    \kq{``훈련 점수가 높다''는 사실은 잘 배웠다는 증거가 아니라,
    \textbf{외웠을 가능성}을 알리는 신호일 수 있으며, 모델이 좋은지
    나쁜지를 결정하는 것은 언제나 \textbf{보지 않은 데이터의
    점수}다.} (엄밀한 처리는 Week 6.)
  \end{block}
\end{frame}
